\uuid{2f5K}
\exo7id{6077}
\titre{exo7 6077}
\auteur{queffelec}
\organisation{exo7}
\datecreate{2011-10-16}
\isIndication{false}
\isCorrection{false}
\chapitre{Continuité, uniforme continuité}
\sousChapitre{Continuité, uniforme continuité}
\cpge{mpsi-mp2i.limites-continuite.b, pcsi.limites-continuite.b, ptsi.limites-continuite.b, mpsi-mp2i.limites-continuite.c, pcsi.limites-continuite.c, ptsi.limites-continuite.c, bcpst1.limites-continuite.c, mp-mpi.topologie-espaces-normes.e, pc.espaces-vectoriels-normes.e, psi.espaces-vectoriels-normes.e, pt.fonctions-deux-trois-variables.b, mp-mpi.topologie-espaces-normes.f}
\module{Topologie}
\niveau{L3}
\difficulte{}

\contenu{
\question{
Soit $C$ l'espace des fonctions continues réelles sur $[0,1]$ muni de la métrique
$d(f,g)=\int_0^1\vert f-g\vert\ dx$, puis de la métrique
$d(f,g)=\sup_x\vert f(x)-g(x)\vert$. Vérifier que l'applica\-tion
$f\to \int_0^1 f dx$ de
$C$ dans
$\Rr$ est continue dans les deux cas.
}
}
